% !TeX root = ../main.tex

\documentclass[../main.tex]{subfiles}
\begin{document}

\section{Kết luận về giải pháp và kết quả đạt được}

\subsection{Tổng quan giải pháp}

Quá trình phân tích và tổng hợp hệ thống điều khiển lưu lượng dầu nhiên liệu và phụ tải lò hơi đã hoàn thành toàn bộ yêu cầu: khảo sát thực trạng công nghệ, xây dựng sơ đồ điều khiển, mô hình hóa, nhận dạng mô hình, tổng hợp bộ điều khiển và kiểm chứng mô phỏng.

Giải pháp được lựa chọn là cấu trúc điều khiển hai vòng (cascade) phân cấp: vòng trong điều khiển lưu lượng dầu qua van FCV để triệt tiêu nhiễu thủy lực ngay tại gốc; vòng ngoài điều khiển áp suất hơi mới của bộ phận sinh hơi. Cấu trúc này tối ưu đáp ứng đặc tính động học khác biệt giữa hai phân hệ (van đáp ứng trong vài giây, lò hơi có quán inertia lớn với thời gian quá độ hàng chục giây) và đáp ứng các yêu cầu khắt khe của hệ thống nhiệt điện công nghiệp.

\subsection{Kết quả mô hình hóa và nhận dạng}

Dựa trên dữ liệu thực nghiệm đã qua làm sạch, hai mô hình quán tính bậc nhất có trễ (QTB1T) và quán tính bậc hai có trễ (QTB2T) được xác định phù hợp cho hai đối tượng điều khiển:

\begin{enumerate}
    \item \textbf{Vòng trong (van điều khiển lưu lượng)}:
    \[
        G_v(s) = \frac{0{,}8836}{4{,}8694\,s + 1}\,e^{-0{,}8213\,s}
    \]
    với hệ số xác định $R^2 = 0{,}9987$, khẳng định mô hình phù hợp gần như hoàn toàn với dữ liệu thực nghiệm.

    \item \textbf{Vòng ngoài (bộ phận sinh hơi)}:
    \[
        O_1(s) = \frac{1{,}2147}{(6{,}0356\,s + 1)^2}\,e^{-4{,}0457\,s}
    \]
    với $R^2 = 0{,}9975$, phản ánh chính xác đặc tính quá độ hình chữ S của quá trình truyền nhiệt hai giai đoạn trong lò hơi.
\end{enumerate}

\subsection{Tổng hợp bộ điều khiển và kiểm chứng}

Áp dụng phương pháp tổng hợp dựa trên mô hình kết hợp IMC, Direct Synthesis và Bền vững tối ưu, các bộ điều khiển của hệ thống hai vòng được tổng hợp với các thông số chính như sau:

\begin{table}[htbp]
\centering
\caption{Thông số bộ điều khiển hai vòng (phương pháp được lựa chọn)}
\label{tab:controller_params}
\begin{tabular}{|l|c|c|c|c|c|c|}
\hline
\textbf{Vòng} & \textbf{Bộ ĐK} & \textbf{Phương pháp} & $K_p$ & $T_i$ (s) & $T_d$ (s) & $\lambda$ (s) \\
\hline
Vòng trong & PI & IMC & $3{,}3532$ & $4{,}8694$ & -- & $0{,}8213$ \\
\hline
Vòng ngoài & PID & Bền vững tối ưu & $1{,}2595$ & $12{,}0712$ & $3{,}0178$ & -- \\
\hline
\end{tabular}
\end{table}

Kết quả kiểm chứng mô phỏng trên Python với thư viện \texttt{scipy.signal} cho thấy hệ thống hoạt động ổn định và đáp ứng toàn bộ chỉ tiêu:

\begin{itemize}
    \item Vòng trong (lưu lượng dầu): thời gian xác lập $T_{s2} \approx 6$\,s, phản ứng cực kỳ nhanh nhạy, triệt tiêu nhiễu thủy lực ngay tại vòng trong trước khi lan truyền.
    \item Vòng ngoài (áp suất hơi): thời gian xác lập $T_{s1} \approx 32$\,s, độ vọt lố $< 5\%$, hầu như không dao động.
    \item Sai lệch tĩnh: $e_{ss} = 0$ (nhờ khâu tích phân I), đảm bảo áp suất hơi luôn bám sát giá trị đặt.
    \item Độ bền vững: hệ thống ổn định khi thông số đối tượng thay đổi trong phạm vi $\pm 20\%$.
\end{itemize}

\subsection{Kết luận và ý nghĩa}

Cấu trúc điều khiển hai vòng đã chứng minh hiệu quả vượt trội so với cấu trúc đơn vòng, đặc biệt phù hợp với đặc tính trễ và quán inertia lớn của lò hơi nhiệt điện. Kết hợp IMC (đơn giản, phù hợp đối tượng QTB1T của van) và Bền vững tối ưu (kiểm soát chặt overshoot, phù hợp đối tượng QTB2T có trễ lớn của lò) là lựa chọn hợp lý, cân bằng tốc độ đáp ứng và độ ổn định.

Về hiệu quả thực tiễn, giải pháp hứa hẹn mang lại lợi ích kinh tế (giảm tiêu thụ dầu nhiên liệu, tiết kiệm chi phí vận hành), xã hội và môi trường (giảm phát thải khí độc CO, NO$_x$, SO$_x$) đồng thời nâng cao độ an toàn nhờ các cơ chế bảo vệ liên động (MFT) được tích hợp trong hệ thống.

Hướng phát triển tiếp theo có thể tập trung vào việc ứng dụng các kỹ thuật điều khiển dự báo (Smith Predictor, MPC) cho vòng ngoài để rút ngắn thời gian quá độ do trễ thuần của lò hơi, hoặc tích hợp thuật toán tự chỉnh định thích nghi (Adaptive Auto-Tuning) để duy trì hiệu suất tối ưu khi đặc tính thiết bị suy giảm theo thời gian.

\end{document}
